\section{Ownership Contract}

RCASE owns child-package composition and a claim matrix. A case bundle is
allowed to say, for each public claim, which child packages are being cited,
what their hashes are, and which child verifier results, caveats, conflicts, or
unsupported checks remain attached to the claim. The planned shapes are
\texttt{claims/claim-boundary.json} and a child-package index, but those are
reserved interfaces rather than implemented schemas.

\begin{center}
\small
\begin{tabularx}{\textwidth}{XX}
\toprule
RCASE may & RCASE must not \\
\midrule
Organize child packages into a public evidence bundle. & Create facts that are
not available in the cited children. \\
Cite child package hashes and external citations. & Replace child hashes with a
case-level assertion that cannot be traced back. \\
Summarize child package claims in a claim matrix. & Loosen caveats, hide
unsupported checks, or make an inconclusive child result look conclusive. \\
Preserve conflicts between child packages. & Merge conflicting claims without
recording the conflict metadata. \\
Order exhibits for a public reader. & Move ownership of counts, lineage,
audits, custody, maps, or certification into the case layer. \\
\bottomrule
\end{tabularx}
\end{center}

The contract is intentionally narrow. RCASE can make a case easier to read and
easier to cite, but the evidence remains owned by the child packages that
declared or derived it.
